%% Introduction (nuclear-application style): motivation + literature + gap + contributions + outline
%% Former Related-work material is merged here; see also sections/01_introduction_zh.md

Nuclear power plants (NPPs) continuously acquire neutron and instrument-channel measurements for
surveillance and condition monitoring.
Ex-core and in-core detectors yield high-volume time series that support power surveillance, sensor
health assessment, and early anomaly detection, yet the resulting records are typically sparsely
labeled.
In practice, day-to-day diagnosis still depends on authoritative technical manuals, including IAEA
guidance on online monitoring (OLM) for instrument channels and for process and component condition
monitoring~\citep{iaea2008olm1,iaea2008olm2}.
Hashemian~\citep{hashemian2011online} reviewed OLM technologies that extend calibration intervals and
strengthen channel performance assessment.
Beyond calibration surveillance, reactor-physics studies treat neutron-noise analysis as a
non-intrusive probe of vibrations, flow anomalies, and core dynamics~\citep{pazsit2010noise,demaziere2022cortex},
while data-driven methods have automated validation of neutron-detector operation and noise-signal
integrity under constrained detector sets~\citep{tambouratzis2020neutron,tambouratzis2020decomposition}.
Deep-learning fault diagnosis for NPPs continues to improve accuracy on labeled accident and fault
scenarios~\citep{elbordany2024fault}; however, such pipelines usually assume fixed feature schemas or
predefined fault libraries.
Exploratory diagnosis of neutron-current CSV streams often lacks gold-standard fault tags, so
operators must still map observed spikes, dropouts, and collinear channel groups to clauses in I\&C
manuals.
Rule-only workflows are brittle across diverse plant conditions, whereas large language models (LLMs)
used alone may produce fluent but ungrounded recommendations.
These observations motivate language-mediated assistants that observe measurement data with weak
priors and then ground interpretations in citable manuals.

Retrieval-augmented generation (RAG)~\citep{lewis2020rag} addresses part of this grounding problem by
conditioning generation on passages retrieved from an external corpus.
In nuclear applications, the corpus is typically regulatory text, plant technical specifications, work
orders, or design and safety manuals rather than open-domain web content.
Idaho National Laboratory studies have shown that RAG can improve diagnostic question answering and
predictive-maintenance explanations for NPP systems, while emphasizing factuality evaluation because
nuclear coverage in pretraining remains sparse~\citep{reeves2024reliability,lin2025cws}.
Work on power-uprate document interpretation further illustrates both the utility of RAG over
thousands of technical pages and recurring failure modes, including incorrect retrieval,
misinterpretation of retrieved text, reliance on parametric knowledge outside the evidence,
hallucination, and refusal~\citep{mapes2025uprate}.
Agentic, citation-enforcing frameworks such as RADIANT-LLM target local, provenance-tracked decision
support for safety-critical nuclear workflows and report marked reductions in hallucination relative
to ungrounded commercial chatbots~\citep{ndum2026radiant}.
Taken together, these studies indicate that manual-grounded RAG is a practical route to more
trustworthy nuclear language assistance.
Nevertheless, most nuclear RAG demonstrations begin from a user question about documents
(procedures, technical specifications, or maintenance records).
Far fewer systems couple retrieval to an upstream analysis of neutron measurement CSV data, where the
query itself must be constructed from observed channel phenomena rather than from free-text operator
prompts.

Tool-using agents provide a complementary mechanism for the measurement side of the pipeline.
ReAct-style agents~\citep{yao2023react} interleave reasoning with tool calls, enabling iterative
observation--action loops over local analyzers and external interfaces.
For neutron-current diagnosis, this paradigm suggests a persistent workspace blackboard together with
specialized probes---profiling, zero-prior statistical discovery, anomaly and regime segmentation, lag
checks, and spatial snapshots---before any manual retrieval is invoked.
On the knowledge side, practical RAG stacks~\citep{liu2023llamaindex} rely on chunking, dense
embedding, and top-$k$ retrieval, while evaluation frameworks such as RAGAs~\citep{es2024ragas}
emphasize faithfulness and context relevance of generated answers.
Nuclear RAG studies already caution that retrieval errors can propagate into unsafe
advice~\citep{mapes2025uprate}, and corpus-scaling experiments show that knowledge-base growth alters
context precision~\citep{ndum2026radiant}.
What remains under-characterized for neutron-manual corpora is the relationship between a compact
authoritative \emph{core} library and a larger same-domain \emph{non-core} supplement: as the
core--non-core ratio changes, cosine similarity scores may rise while core-document purity collapses,
especially under short Chinese versus English queries and under strict chunk identity versus
document-level relevance.
General RAG benchmarks rarely stress bilingual nuclear terminology or core-versus-supplement purity
under controlled library expansion.

Against this background, the present work addresses three coupled challenges:
(1)~how to extract spikes, dropouts, collinearity or redundancy, and regime changes from
neutron-current CSV data with minimal prior assumptions;
(2)~how to align observed phenomena with manual clauses while retaining citable provenance; and
(3)~whether dense retrieval over same-domain technical corpora remains reliable after corpus
expansion, particularly under Chinese versus English query formulations.
The main contributions are as follows.
First, we develop a zero-prior, tool-using diagnostic agent with a workspace blackboard for
neutron-current analysis.
Second, we introduce a report-driven multi-stage RAG protocol that converts condition reports into
structured JSON queries, retrieves passages from IAEA-centered nuclear I\&C manuals, and produces
expert-style diagnoses with explicit evidence identifiers.
Third, we conduct a same-domain retrieval-robustness study that quantifies corpus dilution, language
effects, chunking choices, and the gap between strict chunk matching and document-level relevance.
\Cref{fig:overview} summarizes the end-to-end pipeline.

\begin{figure}[t]
  \centering
  \includegraphics[width=0.98\linewidth]{fig1_overview}
  \caption{End-to-end pipeline: neutron-current CSV is analyzed by a zero-prior tool-using agent
  that maintains a workspace blackboard and produces a condition report; the report is compressed
  into structured JSON to drive manual-grounded RAG with citable diagnoses. A side branch evaluates
  retrieval robustness under same-domain corpus expansion.}
  \label{fig:overview}
\end{figure}

The remainder of this paper is organized as follows.
\Cref{sec:system} presents the data resources and system design, including the diagnostic agent,
manual corpus indexing, and report-driven RAG diagnosis.
\Cref{sec:eval} describes the evaluation protocol for retrieval robustness.
\Cref{sec:results} reports experimental results and discussion.
\Cref{sec:concl} concludes the paper and outlines future work.
